\section{Referencial teórico}

Esta seção organiza os fundamentos teóricos que sustentam a pergunta de pesquisa e a delimitação
conceitual do projeto. O argumento está estruturado em três eixos complementares: a relação entre
inteligência artificial generativa, RAG e recuperação da informação; as especificidades da informação
legislativa como acervo público; e os critérios de avaliação aplicáveis a sistemas RAG em domínios
institucionais.

\subsection{Inteligência artificial generativa, RAG e recuperação da informação}

Os grandes modelos de linguagem ampliaram a capacidade de interação com sistemas computacionais por
meio de linguagem natural, viabilizando tarefas de sumarização, classificação, geração de texto,
resposta a perguntas e apoio à análise documental \parencite{brownEtAl2020fewShotLearners,
zhaoEtAl2023llmSurvey,minaee2024llmSurvey}. As características dominantes desses modelos se apoiam
em avanços associados \pita{à arquitetura \textit{Transformer}}, combinados a estratégias de
pré-treinamento em larga escala e \pita{alinhamento} para seguir instruções
\parencite{vaswani2017attention,brownEtAl2020fewShotLearners,ouyangEtAl2022instructGpt}.
Apesar dessas capacidades, LLMs não são mecanismos de recuperação documental por definição: suas
respostas podem refletir padrões aprendidos no treinamento, inferências probabilísticas e lacunas
de atualização, sem indicar de forma confiável quais documentos sustentam cada afirmação
\parencite{jiEtAl2023hallucinationSurvey,anhHoang2025hallucinations,dahlEtAl2024largeLegalFictions}.

O RAG surge como uma estratégia para aproximar geração textual e recuperação da informação. Na
formulação clássica, a arquitetura combina um mecanismo de recuperação sobre uma base documental
externa com um modelo generativo que produz a resposta a partir do contexto recuperado
\parencite{lewisEtAl2020rag} \pita{(conforme ilustrado na Figura \ref{fig:rag-pipeline})}. Revisões
recentes classificam essas arquiteturas em configurações ingênuas, avançadas e modulares, destacando
que o desempenho depende de todo o \textit{pipeline}, não apenas do modelo final de geração
\parencite{gao2023ragSurvey,gao2024modularRag}. Componentes como extração documental, segmentação em
\textit{chunks}, enriquecimento por metadados, geração de \textit{embeddings}, recuperação vetorial,
reranqueamento \pita{(ou reclassificação)}, seleção de contexto e \textit{prompting} orientado a
evidências influenciam diretamente a qualidade das respostas \parencite{gupta2024comprehensiveRag,
nie2024textEmbeddingSurvey}.

No componente de recuperação, o problema se aproxima da tradição de recuperação da informação, que envolve representação de documentos, formulação de consultas, ranqueamento e avaliação da pertinência dos resultados recuperados \parencite{manningEtAl2008informationRetrieval}. Em sistemas RAG, essa etapa é frequentemente operacionalizada por recuperação densa ou vetorial, em que representações semânticas são usadas para localizar passagens relevantes ao contexto da consulta \parencite{karpukhinEtAl2020dpr,thakurEtAl2021beir}.

Essa literatura é central para \pita{a presente investigação} porque desloca a análise de uma
discussão genérica sobre \pita{o uso de inteligência artificial} para uma questão avaliativa
delimitada: em que condições uma arquitetura RAG recupera evidências pertinentes, gera respostas
verificáveis e reconhece \pita{os} limites de \pita{seu} escopo. \pita{Tal} delimitação exige
distinguir \pita{analiticamente}, no desenho da pesquisa, \pita{três camadas complementares --}
a recuperação documental, a geração de respostas e o julgamento das respostas\pita{ --, de modo a
preservar a articulação entre elas sem confundi-las.}

\subsection{Informação legislativa, transparência e acervos parlamentares}

O acesso à informação pública é dimensão central da administração democrática, mas a disponibilidade formal de documentos não elimina barreiras de compreensão, localização e uso qualificado. No domínio legislativo, discursos, debates e registros parlamentares são fontes relevantes para reconstruir posicionamentos, agendas, conflitos, justificativas e memória institucional \parencite{skubicFiser2024discourseReview}. Ao mesmo tempo, esses acervos podem ser difíceis de explorar por mecanismos convencionais de busca, sobretudo quando o usuário não conhece previamente a terminologia exata, os atores envolvidos ou a organização dos registros \parencite{bandeiraBernardes2016information}.

No caso brasileiro, estudos recentes da Consultoria Legislativa do Senado mostram que os pronunciamentos em plenário formam um \textit{corpus} extenso, dinâmico e analiticamente relevante, com variações de volume, extensão, interatividade e complexidade textual ao longo do tempo \parencite{blanco2025plenarioPalanqueEstudio,blanco2026letrasGraudas}. Em paralelo, pesquisas sobre organização do conhecimento legislativo indicam que o uso de inteligência artificial nesse domínio precisa considerar estrutura informacional, explicabilidade, mediação documental e confiabilidade dos resultados \parencite{martelloteViola2025organizacaoConhecimento}.

Assim, a informação legislativa não deve ser tratada apenas como massa textual a ser processada.
Trata-se de um recurso institucional, marcado por regras de produção, metadados, autoria, temporalidade,
retificações, usos políticos e requisitos de \textit{accountability} \parencite{bandeiraBernardes2016information,
martelloteViola2025organizacaoConhecimento}. Essa característica reforça a necessidade de combinar inovação técnica com critérios de avaliação compatíveis com o uso público e institucional da informação legislativa.

\pita[caixa,azul]{Não é uma observação de correção, mas eu senti falta de mais detalhes sobre o domínio alvo de sua pesquisa,
como o volume de dados, a granularidade dos documentos, a diversidade de tópicos e a eventual complexidade das perguntas que você
pretende responder. Isso ajudaria a contextualizar melhor o desafio de avaliação do sistema RAG no contexto legislativo.}

\subsection{Avaliação de sistemas RAG em domínios institucionais}

A avaliação de sistemas RAG apresenta desafios próprios porque uma resposta final aparentemente correta pode ocultar
recuperação deficiente, uso inadequado das fontes, extrapolação além do contexto ou julgamento automatizado permissivo.
Essa dificuldade se relaciona ao caráter multidimensional da avaliação de sistemas baseados em modelos de linguagem, nos
quais diferentes dimensões de desempenho, confiabilidade e uso precisam ser analisadas de forma articulada
\parencite{liangEtAl2023helm}. \textit{Corpus} anotados em nível de afirmação e palavra, como RAGTruth, tornam
observáveis conflitos e conteúdos sem suporte que podem desaparecer em uma nota agregada \parencite{niuEtAl2024ragtruth}.
Por isso, a literatura recomenda distinguir métricas de recuperação, critérios de avaliação da geração, avaliação humana,
avaliação automatizada baseada em LLMs (\textit{LLM-as-a-judge}) e análise qualitativa de erros \parencite{saadFalconEtAl2024ares,esEtAl2023ragas,
pipitoneAlami2024legalbenchRag}. \textit{Benchmarks} recentes também apontam a importância de testar diferentes famílias
de perguntas, como questões factuais, comparativas, temporais, multietapas, numéricas e não respondíveis
\parencite{zhuEtAl2025rageval}.

Entre \pita{as} técnicas \pita{de avaliação}, a \textit{LLM-as-a-judge} merece atenção específica neste projeto, porque
permite avaliar respostas geradas com base em rubricas explícitas \parencite{liuEtAl2023geval,liuEtAl2025judgeRag}. Estudos
recentes também mostram que a consistência do juiz depende da formulação da avaliação e das dimensões observadas
\parencite{liuEtAl2025judgeRag}. Essa abordagem pode ajudar a escalar avaliações, mas não elimina a necessidade de validação
humana, pois o julgamento automatizado pode variar conforme modelo, rubrica, formato de entrada e complexidade da tarefa.
Em domínios públicos, esse ponto é decisivo: uma avaliação automatizada não deve ser confundida com certificação de qualidade
institucional.

\pita{Considerando essa} literatura, \pita{este trabalho} organiza a avaliação humana em cinco dimensões operacionais:
correção factual da resposta; aderência ao contexto recuperado \pita{(\textit{groundedness})}; suporte e precisão das citações;
completude proporcional à pergunta; e reconhecimento de insuficiência documental. Essa composição é uma síntese operacional,
e não uma rubrica canônica única: ARES e RAGAS sustentam a separação entre contexto e resposta, Trautmann e RAGTruth informam
a fidelidade e a análise de afirmações, Zhang fundamenta o suporte graduado de citações e GaRAGe fundamenta a avaliação de
suficiência e deflexão \parencite{saadFalconEtAl2024ares,esEtAl2023ragas,trautmannEtAl2024groundedness,niuEtAl2024ragtruth,
zhangEtAl2024fineGrainedCitation,sorodocEtAl2025garage}. A dimensão de insuficiência documental é uma adaptação explícita ao
contexto parlamentar e deverá ser calibrada com avaliadores humanos.

Em domínios públicos, diretrizes e estudos sobre IA enfatizam transparência, finalidade pública, supervisão humana e
responsabilização \parencite{ipu2024aiParliaments,wfd2023guidelinesAiParliaments,deAlmeidaSantos2025aiGovernance}. Estudos
institucionais sobre monitoramento parlamentar acrescentam que transparência formal precisa ser acompanhada por mecanismos
de compreensão, participação e \textit{accountability} \parencite{geunis2023parliamentaryMonitoring}. A meta-avaliação de
avaliadores também precisa ser considerada: o \textit{benchmark} GroUSE mostra que uma boa correlação agregada com outro
juiz não assegura cobertura de casos-limite, recomendando testes unitários para verificar contradição, insuficiência e
relevância no protocolo \parencite{mullerEtAl2025grouse}. Como apoio de organização, o \textit{survey} de
\textcite{yuEtAl2024ragEvaluationSurvey} sintetiza as dimensões de recuperação e geração e reforça que a avaliação deve
preservar a articulação entre ambas. Nesta pesquisa, esses elementos são tratados como contexto institucional para justificar a necessidade de avaliação criteriosa, e não como dimensão autônoma do protocolo. O protocolo concentra-se em aspectos diretamente observáveis no sistema RAG: recuperação documental, geração de respostas, indicação das fontes, explicitação de insuficiência documental e verificabilidade das fontes. A possibilidade de recusar quando o contexto é insuficiente deve ser tratada como comportamento avaliável, pois benchmarks recentes mostram que respostas tendem a incorporar conhecimento não ancorado e a evitar a deflexão \parencite{sorodocEtAl2025garage}.

Nesse contexto, o protocolo deve ser entendido como uma organização explícita de procedimentos, métricas, rubricas e critérios de interpretação voltados a avaliar, de forma reprodutível, dimensões complementares de um sistema RAG.

\pita{\subsection{Trabalhos relacionados e protocolos de avaliação de RAG}}

\pita[caixa,vermelho]{Trabalho relacionados deixou de ter uma seção própria, e passou a ser parte do referencial teórico.}

\pita{Esta seção} posiciona a dissertação diante de estudos que combinam recuperação, geração, avaliação de
\pita{\textit{grounding}} e uso de acervos jurídicos ou parlamentares. A \pita{seleção dos estudos} considera
\pita{\textit{corpus}}, tarefa, unidade de avaliação, participação humana e transferência para os discursos
do Senado. \pita{A Tabela~\ref{tab:trabalhos-relacionados} resume os trabalhos selecionados, com foco em suas
contribuições para a avaliação de RAG e sua aplicabilidade ao contexto parlamentar brasileiro.}

\begin{table}[hb!]
\centering
\caption{Comparação de estudos diretamente relacionados ao protocolo.}
\label{tab:trabalhos-relacionados}
\small
\begin{tabular}{p{2.7cm}p{2.8cm}p{3.1cm}p{2.5cm}}
\toprule
Estudo & Corpus/domínio & Foco de avaliação & Transferência para o Senado \\
\midrule
ARES & QA em múltiplos domínios & \textit{Groundedness}, correção e completude & Rubrica geral; requer calibração local \\
LegalBench-RAG & Questões jurídicas em inglês & Recuperação e resposta fundamentada & Alta proximidade jurídica; baixa proximidade institucional \\
\textcite{trautmannEtAl2024groundedness} & QA jurídico & \textit{Groundedness} por similaridade, NLI e LLM & Base para fidelidade ao contexto \\
RAGTruth & Respostas RAG em vários domínios & Alucinação em nível de caso e palavra & Taxonomia de erros; idioma não parlamentar \\
\textcite{zhangEtAl2024fineGrainedCitation} & Texto gerado com citações & Suporte completo, parcial ou ausente & Base para verificabilidade graduada \\
GaRAGe & QA com passagens anotadas & \textit{Grounding} relevante e deflexão & Base para perguntas não respondíveis \\
\textcite{reuterEtAl2025reliableRetrieval} & Grandes bases jurídicas & Erro no nível do documento e \textit{chunking} & Hipótese de erro documental a testar \\
eRAG & Três conjuntos de QA & Utilidade \textit{downstream} de cada documento & Liga recuperação ao desempenho da resposta \\
GroUSE & QA fundamentado controlado & Meta-avaliação por testes unitários & Calibra o \textit{LLM-as-a-judge} \\
\bottomrule
\end{tabular}
\end{table}

\pita[caixa, vermelho]{O TERMO "QA" É UMA SIGLA AINDA NÃO EXPLICADA. PENSO QUE SE REFIRA A "QUESTION ANSWERING", MAS ISSO DEVERIA
SER EXPLICADO NA PRIMEIRA OCORRÊNCIA.}

\pita[caixa,vermelho]{O QUE É NLI? DEFINIR A SIGLA NA PRIMEIRA OCORRÊNCIA.}

\pita[caixa, vermelho]{NESTA TABELA SURGE A PRIMEIRA REFERÊNCIA AO TERMO GROUNDEDNESS, QUE AINDA NÃO FOI DEFINIDO. A DEFINIÇÃO DEVERIA VIR ANTES,
OU O TERMO EM PORTUGUÊS USADO NO REFERENCIAL TEÓRICO ADOTADO. SERIA UMA DAS 5 DIMENSÕES OPERACIONAIS ESPECIFICADAS, PENSO, TALVEZ
"ADERÊNCIA AO CONTEXTO". ISSO VALE PARA TODAS AS OCORRÊNCIAS SEGUINTES DO TERMO GROUNDEDNESS. VALE TAMBÉM PARA O TERMO "GROUNDING" QUE
SURGE MAIS ADIANTE.}

\pita[caixa, vermelho]{DEFINIR O CONCEITO DE UTILIDADE DOWNSTREAM. ENTENDO QUE SERIA A CAPACIDADE DE UM DOCUMENTO RECUPERADO DE MELHORAR O
DESEMPENHO DO MODELO NA TAREFA DE GERAÇÃO, MAS ISSO NÃO ESTÁ CLARO AQUI.}

\pita{A subseção seguinte detalha os sistemas RAG existentes nos domínios jurídico, legislativo e parlamentar, enquanto a
subseção final discute protocolos, métricas e benchmarks de avaliação de RAG, sintetizando as lições aprendidas
e o posicionamento da dissertação em relação a esses trabalhos.}

\subsubsection{Sistemas RAG nos domínios jurídico, legislativo e parlamentar}

Aplicações legislativas e parlamentares existentes cobrem tarefas distintas. KamerRaad e Talk to the NDAA exploram interfaces de
consulta em debates e legislação; RAGAR trata verificação política; LexDrafter trata terminologia legislativa; e LegalBench-RAG
organiza evidência jurídica. Essas fontes demonstram a variedade de tarefas, mas não formam um benchmark comum. LegisSearch e o
sistema do World Avatar acrescentam grafos, metadados e consultas relacionais.~\pita{Esta} dissertação se diferencia por avaliar
um acervo real de discursos do Senado brasileiro, com perguntas e critérios definidos para recuperação, geração e verificabilidade.

\pita{\textcite{reuterEtAl2025reliableRetrieval}} identificam \pita{a ocorrência dos que eles chamam de} \textit{Document-Level Retrieval
Mismatch} em grandes bases legais e mostram que a recuperação pode apontar para o documento errado mesmo quando o trecho parece
semanticamente próximo.~\pita{Os autores propõem uma abordagem de \textit{Summary-Augmented Chunking} (SAC) para reduzir o erro, mas não
testam sua transferência para outros domínios, embora a} observação \pita{seja} diretamente relevante para discursos parlamentares,
nos quais autoria, sessão, data e proposição podem distinguir documentos linguisticamente semelhantes. O resultado\pita{, contudo, }
deve ser tratado como hipótese de erro a \pita{ser testada} no \textit{corpus} do Senado, não como ganho transferível do SAC.

\subsubsection{Protocolos, métricas e benchmarks de avaliação de RAG}

ARES, RAGAS, G-Eval, ConsJudge e RAGEval cobrem dimensões complementares, mas usam tarefas, modelos e formas de validação diferentes.
Trautmann et al. oferecem uma referência mais próxima para a dimensão de \textit{groundedness}: comparam similaridade, NLI e LLMs em perguntas
jurídicas e reportam macro-F1 de 0,80 para o melhor método no corpus estudado \parencite{trautmannEtAl2024groundedness}. O resultado
apoia a inclusão de uma rubrica de aderência ao contexto, preservando a necessidade de validação local.

\textcite{salemiZamani2024erag} propõem o eRAG, que estima a utilidade de cada documento recuperado pelo efeito que produz quando usado individualmente
pelo gerador; essa perspectiva aproxima a métrica de recuperação do desempenho \textit{downstream}. \textcite{niuEtAl2024ragtruth}
criam o RAGTruth, com anotações em nível de caso e palavra, permitindo observar afirmações contraditórias ou sem suporte.
\textcite{zhangEtAl2024fineGrainedCitation} mostram que o suporte de uma citação pode ser completo, parcial ou inexistente, e que
nenhuma métrica é superior em todas as tarefas de correlação, classificação e recuperação.
GaRAGe amplia esse problema para perguntas que exigem \textit{grounding} por passagem e deflexão quando o contexto não basta; seus resultados
indicam dificuldades de atribuição e recusa mesmo em modelos recentes \parencite{sorodocEtAl2025garage}.

A avaliação do avaliador também é objeto de estudo. ConsJudge mostra que a consistência do juiz deve ser examinada em diferentes
combinações de dimensões e prompts \parencite{liuEtAl2025judgeRag}. Essa evidência reforça o desenho multimétodo da dissertação:
métricas automáticas devem ser comparadas com julgamentos humanos \pita{com} análise de erros.

\subsubsection{Síntese comparativa}

\pita[caixa,vermelho]{Fabrício, veja como essa subsubseção se integra com a subseção de lacuna de pesquisa. Ajuste o texto para que
ambas se encaixem em termos de argumentação e não se repitam.}

A literatura converge em três pontos: recuperação e geração devem ser separadas analiticamente; suporte documental precisa ser
verificado em unidade menor que a resposta inteira; e a avaliação automática requer calibração e limites explícitos.
A comparação das fontes centrais também recomenda separar três unidades de erro: documento recuperado, trecho recuperado e
afirmação gerada. Reuter et al. sustentam a primeira distinção, LegalBench-RAG a segunda e RAGTruth e Zhang et al. a terceira,
enquanto GroUSE e ConsJudge impõem cautela na interpretação de avaliadores automáticos. Essa decomposição orienta a matriz de análise
da dissertação e evita que uma nota agregada oculte falhas de recuperação ou de atribuição. Ela diverge quanto à unidade ideal de
avaliação, à definição de relevância e ao equilíbrio entre conhecimento paramétrico e contexto recuperado. Os estudos selecionados
também se concentram \pita{no idioma} inglês, legislação ou domínios gerais.

A contribuição proposta ocupa a interseção ainda pouco coberta entre \textit{corpus} parlamentar brasileiro, recuperação documental,
respostas geradas, citações verificáveis, insuficiência documental e validação humana independente. O quadro comparativo a ser
consolidado deve registrar: corpus e idioma; jurisdição; tipo de pergunta; arquitetura; métricas de recuperação; critérios de
\textit{grounding}; tratamento de respostas não respondíveis; participação humana; e disponibilidade de dados e código. A lacuna é, portanto, uma lacuna de integração e adaptação contextual, não a ausência de técnicas isoladas.

\subsection{Lacuna de pesquisa}

A literatura existente oferece bases técnicas para compreender arquiteturas RAG, avaliação de respostas geradas e recuperação da informação, bem como estudos sobre informação legislativa e parlamentos. Entretanto, os trabalhos mais próximos distribuem-se em frentes ainda pouco integradas. Estudos sobre RAG enfatizam arquitetura, recuperação, geração e avaliação automatizada, mas frequentemente tratam domínios genéricos ou jurídicos mais amplos \parencite{gao2023ragSurvey,saadFalconEtAl2024ares,esEtAl2023ragas,pipitoneAlami2024legalbenchRag}. Pesquisas sobre parlamentos e informação legislativa discutem acesso, transparência, discurso parlamentar e organização do conhecimento, mas nem sempre implementam ou avaliam sistemas conversacionais baseados em RAG \parencite{skubicFiser2024discourseReview,blanco2025plenarioPalanqueEstudio,blanco2026letrasGraudas,martelloteViola2025organizacaoConhecimento}. Aplicações recentes em legislação e política, por sua vez, exploram \textit{fact-checking}, redação legislativa, consulta a textos normativos, interfaces conversacionais e arquiteturas híbridas, mas diferem quanto a \textit{corpus}, língua, jurisdição, tarefa e desenho de avaliação \parencite{khaliq2024ragar,chouhanGertz2024lexdrafter,garrisonEtAl2024talkNdaa,rogiersEtAl2024kamerraad,colomboEtAl2025legisSearch,mosbachEtAl2025worldAvatarParliament}.

Desse modo, a lacuna refere-se à escassez de estudos empíricos que combinem, no mesmo desenho de pesquisa e no contexto brasileiro, um acervo parlamentar real, um \textit{pipeline} RAG reprodutível, avaliação formal da recuperação documental, avaliação das respostas geradas e análise da verificabilidade das fontes. No caso desta pesquisa, essa lacuna será examinada em discursos da 56\textordfeminine{} Legislatura do Senado Federal, com atenção simultânea à recuperação, à geração e à verificabilidade das fontes. O foco do estudo recai sobre a forma como a recuperação documental, a geração de respostas e a verificabilidade das fontes podem ser avaliadas de maneira articulada em consultas legislativas mediadas por RAG.

Essa lacuna é especialmente relevante para a administração pública, pois envolve simultaneamente organização de um protocolo avaliativo, avaliação empírica e uso de acervos institucionais de interesse público. A dissertação proposta busca contribuir para esse espaço ao desenvolver uma investigação com critérios explícitos para avaliação de sistemas RAG em acervos legislativos.
